\section{Điều ba bài báo chưa giải quyết}
\label{sec:ch14-dieu-chua-giai-quyet}

Mục trước khép lại trên một phép sửa được chấm điểm bằng dụng cụ chưa ai kiểm.
Mục này gom cả ba bài lại, nói cái gì trong chương chuyển được sang phần còn
lại của sách, rồi giao phần chưa xong cho
chương~\ref{ch:khoang-trong-nghien-cuu}.

\index{đối tượng của bài báo!ở chương này}%
Trước hết là chỗ lệch đối tượng, thứ mà Phần~IV gặp lần thứ sáu ở đây và lần
này gặp ba lần trong một chương. Ba bài đo ba đại lượng, và không đại lượng nào
là độ trung thực của định nghĩa~\ref{def:ch01-do-trung-thuc}. Bài của
mục~\ref{sec:ch14-mot-phong-thu} đo tỉ lệ lời giải thích còn nêu được thuộc
tính thiên lệch, tức là đo một phương pháp giải thích trên một mô hình dựng
riêng để đánh bại nó. Bài của mục~\ref{sec:ch14-shap-lam-vu-khi} đo tỉ lệ phân
loại sai, một đại lượng về bộ phân loại chứ không về lời giải thích. Bài của
mục~\ref{sec:ch14-sua-bang-nhan-qua} đo khoảng cách tới một tập giá trị tham
chiếu do chính người thí nghiệm dựng, cộng với một chỉ số của
chương~\ref{ch:do-do-trung-thuc}. Chữ \enquote{faithful} không có mặt trong bài
nào, và đó là chỗ vắng mặt mạnh nhất mà sách đã kiểm được: không phải sự im
lặng của một bài mà của cả một tầng trong thang đọc.

\index{độ trung thực!cái chuyển được từ chương này}%
Trong ba bài thì có đúng một thứ chuyển được, và nó không nằm ở kết quả của bài
nào. Mục~\ref{sec:ch09-doi-tuong-khac-nhau} đã liệt ba thứ có thể vượt ranh
giới từ chương~\ref{ch:chi-so-khong-do-do-trung-thuc} sang attribution đặc
trưng, và thứ thứ nhất là cách dựng ground truth từ thiết kế. Thiết lập của
mục~\ref{sec:ch14-phep-dung-gian} là thứ ấy, dựng qua một đường khác và bởi một
người có động cơ khác hẳn: bên tấn công phải biết mô hình của mình dựa vào cái
gì thì mới giấu được, nên câu trả lời đúng có sẵn ở mọi điểm dữ liệu thật.
Chương~\ref{ch:khoang-trong-nghien-cuu} nhận thứ ấy kèm cái giá của nó, đã ghi
ở mục~\ref{sec:ch14-ground-truth-ke-tan-cong}: kết luận rút ra là kết luận về
một mô hình được dựng riêng, không phải về những mô hình mà hai phương pháp kia
thường được đem đi giải thích. Hai thứ còn lại trong danh sách của
mục~\ref{sec:ch09-doi-tuong-khac-nhau} không được chương này chạm tới.

\index{tập can thiệp!chỗ khuyết trong định nghĩa}%
Chỗ chưa giải quyết thứ nhất là chỗ mà cả ba bài đều đứng cạnh và không bài nào
bước vào. Mục~\ref{sec:ch14-ground-truth-ke-tan-cong} đã phát biểu nó: một
phương pháp attribution chỉ trung thực tương đối với tập can thiệp của chính
nó, và \tn{phép dựng giàn}{scaffolding} hoạt động được vì không ai bắt buộc
phải nói tập ấy là tập nào. Không bài nào trong ba bài đề nghị rằng một phương
pháp phải công bố tập can thiệp của mình, dù cả ba đều khai thác hoặc sửa hậu
quả của việc không công bố. Đó là một yêu cầu cụ thể đặt lên một dụng cụ tương
lai, và nó khớp với yêu cầu mà mục~\ref{sec:ch13-tran-noi-gi} rút ra từ một
hướng hoàn toàn khác: cái trần lý thuyết nói dụng cụ ấy phải mô tả chỗ lệch,
còn chương này nói nó phải nêu chỗ lệch được đo dưới những can thiệp nào.

\index{tấn công!chưa ai đo trên mô hình thật}%
Chỗ thứ hai là một phép đo chưa ai chạy, và nó rẻ. Cả ba bài đều làm việc trên
những thiết lập mà ground truth đã biết vì có người dựng ra nó, dù người ấy là
bên tấn công hay là tác giả của các phương trình sinh dữ liệu. Không bài nào hỏi
một mô hình bình thường, đã huấn luyện trên dữ liệu thật, thì phần sai lệch của
lời giải thích nằm ở đâu và lớn tới đâu. Bài của
mục~\ref{sec:ch14-shap-lam-vu-khi} tiến gần nhất mà không biết mình tiến gần,
vì tỉ lệ phân loại sai của nó là một phép đo trên mô hình thật, không cần
ground truth, chấm điểm chính bộ điểm số đã dẫn đường cho phép sửa. Nó chỉ
chứng thực được mức quan trọng chứ không chứng thực được độ trung thực, nhưng
nó là một điều kiện cần đo được, và mục~\ref{sec:ch08-dieu-kien-can} đã cho
thấy điều kiện cần cũng có chỗ dùng.

\index{Causal SHAP!vòng lặp không có bước loại}%
Chỗ thứ ba do chính ba bài nêu ra, và ba danh sách công việc tương lai của
chúng đọc cạnh nhau thì thành một hình. Bài của
mục~\ref{sec:ch14-mot-phong-thu} ghi rằng nó chưa so được phương pháp của mình
với LIME và SHAP trên một thang công bằng, và chưa đo xem phương pháp ấy có
diễn giải được gì hay không. Bài của mục~\ref{sec:ch14-shap-lam-vu-khi} muốn
những mô hình phân loại không dựa quá nặng vào từng điểm ảnh riêng lẻ. Bài của
mục~\ref{sec:ch14-sua-bang-nhan-qua} liệt bốn hướng, ba trong số đó nhằm cải
thiện nửa
\tn{khám phá nhân quả}{causal discovery}: xử lý biến ẩn, thay PC bằng thuật
toán khác khi số đặc trưng lớn, và gộp nhiều đồ thị khả dĩ. Hướng thứ tư là xấp
xỉ nhanh hơn cho dữ liệu nhiều chiều, tức nhằm vào chi phí tính giá trị
Shapley~\autocite{p29causalshap}. Không hướng nào trong bốn hướng hỏi làm sao
biết bộ điểm số thu được là đúng.

\index{khoảng trống nghiên cứu!bảy chương Phần IV}%
Với ba danh sách ấy thì sổ ghi các phát biểu hạn chế của Phần~IV đủ bảy chương,
và cái nó cho thấy là hai nhóm chứ không phải một. Năm bài, từ bài 21 tới bài
25, nêu cùng một thứ đang thiếu là một dụng cụ đo độ trung thực đã kiểm định,
và nêu từ năm hướng: đo thấy nó hỏng, ghi nhận rằng không có chuẩn dùng nó, gọi
tên bước tiếp theo chưa ai làm, đếm xem phép đối chiếu với người hiếm tới đâu,
và chạy phép đối chiếu ấy một lần rồi thấy nó không khớp. Ba bài không nêu, và
ba chỗ trống của chúng cũng không giống nhau: bài của
chương~\ref{ch:gioi-han-ly-thuyet} muốn một đại lượng tính được, bài của
mục~\ref{sec:ch14-shap-lam-vu-khi} muốn mô hình tốt hơn, bài của
mục~\ref{sec:ch14-sua-bang-nhan-qua} muốn khám phá nhân quả tốt hơn.

Bài của mục~\ref{sec:ch14-mot-phong-thu} là bài duy nhất trong ba bài của
chương có giao với năm bài kia, và chỗ giao của nó là chỗ mạnh nhất trong cả
sổ. Mục~\ref{sec:ch11-dieu-chua-giai-quyet} đã ghi lần đầu tiên một bài phương
pháp tự nói ra rằng phép so nghiêm túc sẽ cần một phép đo độ trung thực mà nó
không làm. Ở đây điều đó lặp lại từ bên trong một bài có mục đích sửa chữa lời
giải thích, và nói thẳng hơn: các tác giả ghi rằng họ chưa đo xem công cụ của
mình có diễn giải được gì không~\autocite{p27shlime}. Một bài đề xuất cách sửa
lời giải thích, tự nhận chưa kiểm xem lời giải thích sau khi sửa có tốt hơn hay
không, là vòng lặp không có bước loại của
mục~\ref{sec:ch11-dieu-chua-giai-quyet} quay đúng một vòng nữa ở nhánh cuối
cùng còn lại.

\index{giải thích hậu nghiệm!chỗ yếu chung}%
Còn một điều chương này để lại cho cả Phần~V chứ không riêng cho một chương, và
nó là điều mà ba bài cộng lại mới nói được. Cả ba phương pháp bị tấn công, dùng
để tấn công, hay được đem đi sửa đều là phương pháp hậu nghiệm: mô hình huấn
luyện xong trước, rồi một công cụ đứng bên ngoài hỏi nó. Chỗ mà cả ba bài xoay
quanh là chỗ những câu hỏi ấy phải hỏi ở đâu, và câu trả lời luôn là ở những
đầu vào đã bị thay hoặc bị che, tức những đầu vào không phải dữ liệu thật. Phép
dựng giàn biến chỗ ấy thành một cánh cửa, phép tấn công bằng SHAP đi qua nó
theo chiều ngược, và phép sửa bằng cấu trúc nhân quả cố thu hẹp nó lại mà vẫn
phải đứng bên ngoài mô hình. Chương~\ref{ch:nut-that-khai-niem} mở Phần~V bằng
cách bỏ hẳn tư thế đứng bên ngoài ấy: mô hình được dựng sao cho các đại lượng
bên trong nó vốn đã là những khái niệm đọc được, nên không còn phải hỏi nó ở
đâu cả. Phần~V đọc xem cái giá của lối thoát ấy là gì.
